\chapter{先导课程}
\chaptermark{先导课程}
\section{写在前面}
解析几何是高中数学的重要学习内容，在高考中分值占比较高。

不少教辅会把这一部分直接叫作“圆锥曲线”，因为圆锥曲线确实是解析几何中的核心内容。本章除圆锥曲线外，还包括直线与圆的基础内容，先从直线讲起。

解析几何的计算通常从对象和约束入手，把几何条件写成方程，再通过代数变形得到结论。关系较直接时，失误往往出在运算和变形上，因此应选用结构清楚、便于核对的计算方法。
\section{轮换，对称}
先给出几个基本概念，并列出一组常见的式子：
\begin{example}
    观察以下代数式，并归纳它们的共同特点：\vspace{-10pt}
    \begin{gather*}
        abc\\
        a+b+c \\
        ab+bc+ca \\
        a^2+b^2+c^2 \\
        (a+b)(b+c)(c+a)\\
        a^2b + a^2c + b^2a + b^2c + c^2a + c^2b
    \end{gather*}
\end{example}
\begin{solution}
    这些式子的共同点，是交换任意两个变量后，式子都不变。
    
    这个共同点可以直接表述为：若式中出现某一类项，那么交换字母后得到的同类项也都出现。比如在$a+b+c$中，三个字母的出现方式完全一样；在$a^2b + a^2c + b^2a + b^2c + c^2a + c^2b$中，若出现$a^2b$，也会相应出现$a^2c,b^2a,b^2c,c^2a,c^2b$。同类项是成组出现的。
    
    这些式子在交换任意两个变量后都不变，这类式子称为对称式。
\end{solution}
\begin{definition}[对称式]
    对于一个 $n$ 元多项式 $P(x_1, x_2, \dots, x_n)$，若任意排列变量后都有
    \[
    P(x_{i_1}, x_{i_2}, \dots, x_{i_n}) = P(x_1, x_2, \dots, x_n),
    \]
    则称 $P(x_1, x_2, \dots, x_n)$ 为对称式。
\end{definition}
“对称”的意思是各变量在式中的作用完全相同。只要式子中包含某个由特定字母组成的项（例如$a^2b$），那么它一定包含由所有其他字母以同样方式组成的项（即$a^2c,b^2a,b^2c,c^2a,c^2b$）。

判断对称式时，可以直接检查变量交换后表达式是否保持不变。像“$a^2b + a^2c + b^2a + b^2c + c^2a + c^2b$”这样的式子较长，通常改用更紧凑的记号：
\begin{definition}[循环和]
    对三个字母 $a,b,c$，若某一项按
    \[
    a\to b\to c\to a
    \]
    依次轮换后相加，则称所得和为循环和，常记作
    \[
    \sum_{\mathrm{cyc}} f(a,b,c)=f(a,b,c)+f(b,c,a)+f(c,a,b).
    \]
    例如
    \[
    \sum_{\mathrm{cyc}}a^2b=a^2b+b^2c+c^2a.
    \]
\end{definition}

\begin{property}[对称式的性质]
    三元对称多项式常用
    \[
    a+b+c,\qquad ab+bc+ca,\qquad abc
    \]
    作基本表达。若一个式子对 $a,b,c$ 对称，整理时把它化成这三类量的组合，变量地位会更一致，后续代换也更稳。
\end{property}

再列一组式子：
\begin{example}
    观察以下代数式，并归纳它们的共同特点：\vspace{-10pt}
    \begin{gather*}
        a^2b+b^2c+c^2a\\
        a^3b+b^3c+c^3a\\
        \dfrac{a}{b+c}+\dfrac{b}{c+a}+\dfrac{c}{a+b}
    \end{gather*}
\end{example}
\begin{solution}
    这几类式子和前面的对称式不同。交换两个字母后，结果一般会变。比如$a^2b+b^2c+c^2a$中，若把$a$和$b$交换，就会出现原式中没有的$b^2a$项。

    它们也有固定的规律。以$a^3b+b^3c+c^3a$为例，$a^3,b^3,c^3$都会出现，但它们后面跟着的字母是按$a\to b\to c\to a$的顺序轮换的，而不是像对称式那样把所有可能的组合都写出来。
\end{solution}
\begin{example}
    将$(a+b)(b+c)(c+a)$展开，并尽量把结果写成由$a,b,c$对称出现的形式：
\end{example}
\begin{solution}
    展开后整理得：\vspace{-10pt}
    \begin{align*}
    (a+b)(b+c)(c+a)&=(b^2+ac+ab+bc)(a+c)\\
    &=ab^2+b^2c+a^2c+ac^2+a^2b+abc+abc+bc^2\\
    &=(a+b+c)(ab+bc+ca)-abc
    \end{align*}
\end{solution}

